\section{物語プロット}
\subsection{シーン１：森}
うっそうとした森の中。逃げるエルフの少女の息遣い。
背後から下卑た笑い声を上げる賊たちに追い詰められ、少女は木の根に躓いて転ぶ。
賊が手を伸ばしたその時、重苦しい金属音が森に響く。\\
賊の前に立ち塞がったのは、全身錆だらけの騎士。動きは遅く、関節からはギチギチと不快な音がする。
賊は「なんだお前は！」と斬りかかるが、騎士は無言のまま避けもせず、重厚な剣の一撃で賊の武器ごと殴り飛ばす。
華麗な剣技ではない。ただ硬く、重いだけの暴力。賊たちは恐怖して逃げ出す。\\
騎士は少女を一瞥もせず、軋む足取りで去ろうとする。

\subsection{シーン２：街道}
少女は騎士の後を追う。
騎士は立ち止まらず、兜越しのくぐもった声で拒絶する。
\\「ついてくるな。俺はもうガタがきてる鉄屑だ。……子守をしてやる余裕はない。」\\
それでも少女は、騎士の錆びついた足に合わせて一生懸命歩く。
道端に咲く、小さな青い花を見つける少女。
\\「あ！ おじちゃん、これ！」\\
少女は騎士の前に回り込み、背伸びをして、騎士の錆びた腰甲冑の隙間にその花を無理やりねじ込む。
\\「お守り！ 助けてくれたお礼！」\\
無骨で汚れた鉄の塊に、不釣り合いな可憐な花が一輪。
騎士は溜息をつく。それを引き抜こうとするが、ガントレットの指が太すぎて（あるいは、老いで指が震えて）うまく摘めない。
\\「……勝手にしろ。」\\
騎士は諦めて歩き出す。少女は「やった！」と笑顔でその後ろをついていく。

\subsection{シーン３：夜の焚き火}
焚き火のパチパチという音。\\
騎士は兜を被ったまま、固い干し肉を齧っている（隙間から押し込んでいる）。
少女も渡された干し肉を齧ろうとするが、固すぎて噛み切れず、困った顔をしている。
それを見た騎士は、ため息をつくと、無言で少女の持っていた固い肉をひったくり、ナイフで細かく削いで少女に差し出す。
少女はパァッと顔を輝かせ、「ありがとう！」と受け取る。
\\「おじちゃんは、どこへ行くの？」
\\「……終わる場所だ。」
\\「終わる場所？お家？」
\\「……お前たちエルフにはわからんだろうな。」\\
騎士は自分の手を見る。ガントレットの下の手が、小刻みに震えているのを隠すように拳を握る。
\\「俺の時間はもう無い。……何千年も生きるお前とは、見ている世界が違いすぎる。」\\
騎士は焚き火越しに少女を見る。炎に照らされた彼女の顔は、あまりに生気に満ちていて、眩しい。
\\「長く戦いすぎたんだ。……近くにいるだけで、お前にまで錆（死の臭い）が移っちまう。」
\\（これ以上、懐かれるわけにはいかない。俺が死んだ時、この子が悲しむだけだ。）\\
騎士はわざと、ドスの利いた低い声で突き放す。
\\「迷惑なんだよ。明日、エルフの里の近くで別れる。……二度と俺の前に現れるな。」\\
少女は悲しそうに俯くが、何も言い返さない。腰の青い花だけが、風に揺れている。\\

夜が更ける。少女は焚き火のそばで丸くなって眠っている。騎士は座ったまま動かない。
\\「……ん……おじ……ちゃん……」\\
少女が小さく寝言を漏らす。不安そうな、しかしどこか甘えるような声。
騎士はその声にピクリと反応し、ガントレットの拳をギリッと握りしめる。
\\「……すまんな」\\
騎士は音もなく立ち上がり、少女のズレた毛布を首元までかけ直す。

\subsection{シーン４：里の近く}
翌朝。二人は無言のまま歩いている。\\
昨夜の拒絶が重くのしかかり、少女は騎士から数歩遅れて、うつむきながらついてくる。
騎士も振り返らない。これでいい、と自分に言い聞かせるように、ただ前だけを見ている。\\
木漏れ日が差し込み、森が開ける。
\\「……行け。ここまでだ。」\\
騎士が足を止め、冷たく告げる。
少女はハッとして顔を上げ、何か言いたげに口を開く。
\\「……っ、でも……！」\\
しかし騎士の頑なな背中を見て、言葉を飲み込み、唇を噛みしめて俯く。
その時、下卑た笑い声が静寂を切り裂く。

\subsection{シーン５：襲撃}
エルフの少女を取り返しに、再びあの賊たちが現れる。
\\「見つけたぞ、あのエルフだ！」
\\「手間かけさせやがって……。おい、傷物にするなよ？ こいつは金貨の山なんだ。」\\
騎士は即座に少女を背に庇い、抜剣する。老体に鞭を打ち応戦するが、賊の数は前回よりも多い。
\\（死ぬのは怖くないはずだった。だが、今ここで倒れる姿を……この子の未来を閉ざすわけにはいかない！）\\
防戦一方となり、敵の重い一撃が騎士の頭部を捉える。
衝撃音と共にバイザーが砕け散り、騎士はその勢いでガクリと地面に膝をつく。
\\「あっ……！」\\
少女が息を呑む。砕けた仮面の下から、初めて騎士の素顔（片目）が露わになる。
賊の一人が苛立ちながら声を荒げる。
\\「チッ、しぶといジジイだ。……おい、死にたくなきゃ失せろ。俺たちの狙いはガキだけだ。見逃してやる。」\\
しかし騎士は、その言葉など聞こえていないかのように、露出した瞳で彼女を真っ直ぐに見据え、今までで一番大きな、そして優しい声で叫ぶ。
\\「行け！ お前には……未来があるんだ！」\\
その形相と迫力、そして込められた想いに打たれ、少女は一瞬迷った後、唇を噛み締めて踵を返し、全速力で走り出す。

\subsection{シーン６：決戦}
遠ざかる小さな足音を、騎士は背中で聞く。
\\「……そうだ、それでいい。」\\
騎士はよろりと、しかし力強く再び立ち上がる。
無視された賊が舌打ちをする。
\\「あ？ 死にたいのかテメェ」\\
騎士はフッと笑い、剣を構え直して立ちはだかる。
もはや孤独に死ぬためではない。少女の記憶の中に「守られた」という温かい事実を残すために……！
\\「安心しろ。……もう追わせん！」\\
逃げていく少女の背中へ向けた慈愛に満ちた瞳が、賊に向き直った瞬間、戦士の鋭い眼光へと変わる。
騎士は最後の力を振り絞り、賊の群れを迎え討つ。激しい剣戟の中で、腰に挿していた一輪の花が、花弁となって舞い散る。

\subsection{シーン７：別れ}
戦いは終わった。賊は全滅し、騎士もまた、大樹の根元に座り込んだまま動かない。
鎧はあちこちが砕け、体からは力が抜けている。
\\「おじちゃん……っ！」\\
戻ってきた少女が駆け寄り、騎士の冷たい鎧に抱きつく。
騎士は、うっすらと目を開け、困ったように眉を寄せる。
\\「……馬鹿野郎が。……逃げろと、言っただろう……」\\
その時、彼女は気づく。騎士の腰の甲冑の隙間には、花弁の散ってしまった茎だけが、まだしがみつくように残っていることを。
騎士は、砕けたバイザーの隙間から見える目で、泣きじゃくる少女を優しく見つめる。
\\「……泣くな。お前にとっちゃ、瞬き一回分の時間だろうが……」\\
震える手（ガントレット）を上げ、エルフの頬に触れようとする。
\\「お前を守れた。……それで、十分だ……」\\
（スローモーション）少女は涙を流しながら、その手を両手で包み込もうとする。
しかし、あと数センチのところで、ガントレットの重さに耐えかねて、手はガクリと力なく落ちる。
地面に手が落ちた「ゴトッ」という鈍い音だけが響く。\\
森に風が吹き抜け、木々がざわめく音だけが、静寂の中に響き渡る。

\subsection{エンディング（スタッフロール）}
黒画面（または思い出のカットのフラッシュバック）に、スタッフロールが流れる。
観客に「騎士の死」を噛み締めさせる時間。

\subsection{シーン８：数百年後}
季節は巡り、風景は移り変わっていく。
かつて騎士が座り込んでいた場所には、錆びて赤茶色になった兜（バイザーが欠けたもの）だけが、地面に半ば埋もれるように遺っている。
しかし、その場所は寂しくはない。兜を取り囲むように、あの日と同じ種類の「青い花」が一面に咲き乱れている。
\\「……おじさん。今年も……また、春が来たよ。」\\
美しい大人の女性に成長したエルフが、大樹の根元で兜の横に座り、穏やかに語りかけている。
その場所だけが、まるで永遠の春が続いているかのように、光と花に包まれていた。

\textbf{「完」}